\section{Introduction and Examples}

\exercise{1.1}{A cube of i}

$i^3=i^2i=-i$, so the answer is $0-i$.

\exercise{1.2}{A fourth power of i}

$i^4=(i^2)^2=(-1)^2=1$, or $1+0i$.

\exercise{1.3}{A high power of i}

Since $23=4\cdot5+3$ and $i^4=1$, $i^{23}=(i^4)^5i^3=-i$.

\exercise{1.4}{A power of minus i}

$(-i)^{35}=(-1)^{35}i^{35}=-i^3=i$, because $35\equiv3\pmod4$.

\exercise{1.5}{Multiplying two binomials}

Distribute and use $i^2=-1$: $(4-i)(5+3i)=20+12i-5i-3i^2=23+7i$.

\exercise{1.6}{Another binomial product}

$(8+2i)(3-i)=24-8i+6i-2i^2=26-2i$. The real contribution from $-2i^2$ is $+2$.

\exercise{1.7}{A product plus a complex number}

$(2-3i)(4+i)=8+2i-12i-3i^2=11-10i$. Adding $6-5i$ yields $17-15i$.

\exercise{1.8}{Cubing one plus i}

First $(1+i)^2=1+2i+i^2=2i$. Thus $(1+i)^3=2i(1+i)=-2+2i$.

\exercise{1.9}{A fifth power by the binomial theorem}

Expand $(1-i)^5=1-5i+10i^2-10i^3+5i^4-i^5$. Replacing $i^2=-1,i^3=-i,i^4=1,i^5=i$ gives $(1-10+5)+(-5+10-1)i=-4+4i$.

\exercise{1.10}{The modulus of three minus four i}

For $a+bi$ the modulus is $\sqrt{a^2+b^2}$. Hence $|3-4i|=\sqrt{9+16}=5$.

\exercise{1.11}{The modulus of six plus four i}

$|6+4i|=\sqrt{36+16}=\sqrt{52}=2\sqrt{13}$. The modulus is nonnegative.

\exercise{1.12}{Polar form: Three minus four i}

Compute $r=|3-4i|=5$ and divide the real and imaginary coordinates by $r$. This gives
\[3-4i=5\left(3/5+\left(-4/5\right)i\right).
\]
The two displayed normalized coordinates have squared sum $1$, so the parenthesized complex number lies on the unit circle. This is the requested form $r(p+qi)$.

\exercise{1.13}{Polar form: Minus one plus i}

Compute $r=|-1+i|=\sqrt2$ and divide the real and imaginary coordinates by $r$. This gives
\[-1+i=\sqrt2\left(-1/\sqrt2+\left(1/\sqrt2\right)i\right).
\]
The two displayed normalized coordinates have squared sum $1$, so the parenthesized complex number lies on the unit circle. This is the requested form $r(p+qi)$.

\exercise{1.14}{Polar form: Twelve plus five i}

Compute $r=|12+5i|=13$ and divide the real and imaginary coordinates by $r$. This gives
\[12+5i=13\left(12/13+\left(5/13\right)i\right).
\]
The two displayed normalized coordinates have squared sum $1$, so the parenthesized complex number lies on the unit circle. This is the requested form $r(p+qi)$.

\exercise{1.15}{Polar form: Minus three plus five i}

Compute $r=|-3+5i|=\sqrt{34}$ and divide the real and imaginary coordinates by $r$. This gives
\[-3+5i=\sqrt{34}\left(-3/\sqrt{34}+\left(5/\sqrt{34}\right)i\right).
\]
The two displayed normalized coordinates have squared sum $1$, so the parenthesized complex number lies on the unit circle. This is the requested form $r(p+qi)$.

\exercise{1.16}{Fourth roots of one}

Write $z=re^{i\theta}$. Then $r^4=1$, so $r=1$, and $4\theta=2\pi k$. Distinct roots occur for $k=0,1,2,3$:
\[z=e^{i\pi k/2}\in\{1,i,-1,-i\}.\]
Larger integers repeat these roots; a fourth-degree polynomial has no more than four distinct roots.

\exercise{1.17}{Fourth roots of minus one}

Since $-1=e^{i(\pi+2\pi k)}$, the roots are
\[z=e^{i(\pi/4+k\pi/2)},\quad k=0,1,2,3.
\]
In Cartesian form these are $(1+i)/\sqrt2$, $(-1+i)/\sqrt2$, $(-1-i)/\sqrt2$, and $(1-i)/\sqrt2$. Each has modulus $1$ and fourth power $-1$.

\exercise{1.18}{Cube roots of minus eight}

Here $-8=8e^{i(\pi+2\pi k)}$. Thus $r=2$ and $\theta=\pi/3+2\pi k/3$. For $k=0,1,2$ the roots are $1+i\sqrt3$, $-2$, and $1-i\sqrt3$. Distinct arguments modulo $2\pi$ give all three solutions.

\exercise{1.19}{Cube roots of minus twenty-seven i}

Use $-27i=27e^{i(-\pi/2+2\pi k)}$. The roots are $3e^{i(-\pi/6+2\pi k/3)}$ for $k=0,1,2$, namely
\[\frac{3\sqrt3}{2}-\frac32i,\qquad3i,\qquad-\frac{3\sqrt3}{2}-\frac32i.
\]
Cubing triples each argument and cubes the modulus, giving $-27i$ in every case.

\exercise{1.20}{Sixth roots of one}

The six roots are $e^{i\pi k/3}$, $k=0,\ldots,5$. Explicitly:
\[1,\quad\tfrac12+\tfrac{\sqrt3}{2}i,\quad-\tfrac12+\tfrac{\sqrt3}{2}i,\quad-1,\quad-\tfrac12-\tfrac{\sqrt3}{2}i,\quad\tfrac12-\tfrac{\sqrt3}{2}i.
\]
They follow from $r^6=1$ and $6\theta\in2\pi\Z$.

\exercise{1.21}{Sixth roots of minus sixty-four}

Here $r^6=64$ gives $r=2$ and $6\theta=\pi+2\pi k$. Thus $z=2e^{i(\pi/6+k\pi/3)}$, $k=0,\ldots,5$, or
\[\sqrt3+i,\quad2i,\quad-\sqrt3+i,\quad-\sqrt3-i,\quad-2i,\quad\sqrt3-i.
\]
These have six distinct arguments modulo $2\pi$.

\exercise{1.22}{Addition modulo seventeen}

$10+16=26=17+9$, so $10+_{17}16=9$. In modular addition, choose the unique representative between zero inclusive and the modulus exclusive.

\exercise{1.23}{Addition modulo ten}

$8+6=14=10+4$, so $8+_{10}6=4$. In modular addition, choose the unique representative between zero inclusive and the modulus exclusive.

\exercise{1.24}{Real modular addition}

$20.5+19.3=39.8=25+14.8$, so the representative in $[0,25)$ is $14.8$. In modular addition, choose the unique representative between zero inclusive and the modulus exclusive.

\exercise{1.25}{Addition modulo one}

$1/2+7/8=11/8=1+3/8$, hence $1/2+_1 7/8=3/8$. In modular addition, choose the unique representative between zero inclusive and the modulus exclusive.

\exercise{1.26}{Angles modulo two pi}

$3\pi/4+3\pi/2=9\pi/4=2\pi+\pi/4$, so the result is $\pi/4$. In modular addition, choose the unique representative between zero inclusive and the modulus exclusive.

\exercise{1.27}{An irrational modulus}

$2\sqrt2+3\sqrt2=5\sqrt2$, and $\sqrt{32}=4\sqrt2$. Subtracting that modulus once leaves $\sqrt2\in[0,4\sqrt2)$, the answer. In modular addition, choose the unique representative between zero inclusive and the modulus exclusive.

\exercise{1.28}{Inputs outside the underlying set}

The operation $+_6$ on $\R_6=[0,6)$ accepts only ordered pairs of elements of that set. The number $8$ is not an element of $\R_6$, so $5+_6 8$ is not an operation on the stated domain. One could first replace $8$ by $2$, but that is an additional reduction, not the given expression as written.

\exercise{1.29}{A linear equation modulo fifteen}

The congruence is $x+7\equiv3\pmod{15}$, so $x\equiv-4\equiv11\pmod{15}$. The only representative in $\Z_{15}$ is $11$, and $11+_{15}7=3$.

\exercise{1.30}{An angular equation}

Subtract $3\pi/2$ from $3\pi/4$: $x\equiv-3\pi/4\pmod{2\pi}$. Adding $2\pi$ gives $x=5\pi/4\in[0,2\pi)$. Its sum with $3\pi/2$ is $11\pi/4$, whose remainder is $3\pi/4$.

\exercise{1.31}{Twice a residue modulo seven}

We need $2x\equiv3\pmod7$. Since $2\cdot4\equiv1$, multiplication by $4$ gives $x\equiv12\equiv5$. Checking, $5+_7 5=3$. The inverse of $2$ makes the answer unique.

\exercise{1.32}{Three times a residue modulo seven}

The equation is $3x\equiv5\pmod7$. Since $3\cdot5\equiv1$, we get $x\equiv25\equiv4$. Indeed $4+4+4=12\equiv5\pmod7$.

\exercise{1.33}{Twice a residue modulo twelve}

The equation $2x\equiv2\pmod{12}$ means $12\mid2(x-1)$, equivalently $6\mid x-1$. Representatives in $\Z_{12}$ are therefore $x=1$ and $x=7$. Both give remainder $2$ when doubled. Dividing a congruence here also divides its modulus.

\exercise{1.34}{Four times a residue modulo four}

Every residue $x\in\{0,1,2,3\}$ satisfies $4x\equiv0\pmod4$. Hence all four residues are solutions; the coefficient has no multiplicative inverse modulo $4$.

\exercise{1.35}{An isomorphism from eighth roots}

Multiplication of powers corresponds to addition, so the image of $\zeta^m$ must be $5m\pmod8$. For the six requested exponents:
\[\begin{array}{c|rrrrrr}m&0&3&4&5&6&7\\\hline 5m\bmod8&0&7&4&1&6&3\end{array}\]
This also gives $m=2\mapsto2$, as required in the question. Since $\gcd(5,8)=1$, all eight images are distinct.

\exercise{1.36}{An isomorphism from seventh roots}

The image of $\zeta^m$ is $4m\pmod7$ because powers become repeated sums. Thus
\[\begin{array}{c|rrrrrr}m&0&2&3&4&5&6\\\hline4m\bmod7&0&1&5&2&6&3\end{array}\]
Multiplication by $4$ is a permutation of $\Z_7$, since $4$ is relatively prime to $7$.

\exercise{1.37}{Why the proposed image is impossible}

The element $\zeta=e^{i\pi/3}$ generates all six roots, so its image under an isomorphism must generate $\Z_6$. But $4$ generates only $\{0,4,2\}$, since $3\cdot4\equiv0\pmod6$. Alternatively, $\zeta^3=-1\ne1$ would map to $0$, violating injectivity. Thus no isomorphism can have this image.

\exercise{1.38}{Angle addition from Euler formula}

Compute the product in two ways:
\[e^{ia}e^{ib}=e^{i(a+b)}=\cos(a+b)+i\sin(a+b),\]
whereas direct multiplication gives
\[(\cos a+i\sin a)(\cos b+i\sin b)=(\cos a\cos b-\sin a\sin b)+i(\sin a\cos b+\cos a\sin b).
\]
Equality of real parts proves the cosine identity; equality of imaginary parts proves the sine identity. This argument explains why multiplying unit complex numbers adds their angles.

\exercise{1.39}{Multiplying polar forms}

Put $r_j=|z_j|$. Distributing the factors gives
\[z_1z_2=r_1r_2\{(\cos\theta_1\cos\theta_2-\sin\theta_1\sin\theta_2)+i(\sin\theta_1\cos\theta_2+\cos\theta_1\sin\theta_2)\}.
\]
The two addition identities replace the braces by $\cos(\theta_1+\theta_2)+i\sin(\theta_1+\theta_2)$. Thus the moduli multiply and the angles add. If one number is zero, the same product formula is zero regardless of the chosen angle.

\exercise{1.40}{The triple-angle identity}

\begin{enumerate}[label=(\alph*)]
\item Cubing Euler's formula gives $e^{3i\theta}=(\cos\theta+i\sin\theta)^3$. Its real part is $\cos^3\theta-3\cos\theta\sin^2\theta$, since the terms with odd powers of $i$ are imaginary and $i^2=-1$. Hence $\cos3\theta=\cos^3\theta-3\cos\theta\sin^2\theta$.
\item Substitute $\sin^2\theta=1-\cos^2\theta$ to obtain $\cos3\theta=\cos^3\theta-3\cos\theta(1-\cos^2\theta)=4\cos^3\theta-3\cos\theta$.
\end{enumerate}

\exercise{1.41}{Euler's formula from power series}
Start with $e^z=\sum_{n=0}^\infty z^n/n!$ and substitute $z=i\theta$. Separate even and odd powers using $i^{2k}=(-1)^k$ and $i^{2k+1}=i(-1)^k$:
\[
\begin{aligned}
e^{i\theta}
&=\sum_{k=0}^\infty\frac{(-1)^k\theta^{2k}}{(2k)!}
+i\sum_{k=0}^\infty\frac{(-1)^k\theta^{2k+1}}{(2k+1)!}\\
&=\cos\theta+i\sin\theta.
\end{aligned}
\]
The two real series are exactly the given cosine and sine expansions. This is the requested formal derivation; their absolute convergence for every real $\theta$ also justifies separating the two subseries analytically.
